The model is an XGBoost classifier \citep{chen2016xgboost}.
Hyperparameters are fixed across experiments.
Training uses early stopping on the validation split.

Appendix Table \ref{tab:xgb_params} lists key model hyperparameters.
The model uses objective \texttt{binary:logistic}.
Training uses early stopping rounds $=50$ with evaluation metrics AUC and AUC PR on the validation split.

This paper compares three feature configurations.
The first configuration is a time aware target encoding baseline.
The second configuration adds time aggregation features to this baseline.
The third configuration removes target encoding and uses only time aggregation features as a sensitivity check.
Target encoding follows standard practice for high cardinality categorical features \citep{micci2001preprocessing}.

The time aware target encoding baseline is computed at hour resolution with a one hour shift.
For a row at hour $h$, statistics use only labels from hours $< h$.
Appendix Section \ref{sec:appendix_te_details} provides the full definition and implementation details.

Time aggregation features are entity history summaries computed under the no lookahead constraint.
Time is represented at hour resolution.
For an impression at hour $H$, all history based features use only events from hours strictly before $H$.
This excludes same hour information.

The following entity keys are used: device ip, device id, app id, and site id.
For each key and each window, the features include impression counts and smoothed click rates.
For impressions $I$ and clicks $C$ in a window, the impression feature is $\log(1+I)$.
The click rate feature is $(C + \alpha) / (I + \alpha + \beta)$ with $\alpha=1$ and $\beta=10$.

This paper varies two aspects of window design.
The first aspect is the length tuple, expressed in hours back.
The second aspect is the window shape.
The window shapes include trailing windows, a one hour gap variant, bucketized recency buckets, calendar aligned windows, and an event count window defined as the last fifty impressions.

This paper defines a window specification as a length tuple and a shape.
For example, the trailing specification with length tuple $(1, 6, 24, 48, 168)$ uses trailing windows at those hour lengths.

\subsection{Window length tuples}
Table \ref{tab:length_sets} lists the benchmarked length tuples.
Each tuple is evaluated under multiple window shapes.

\begin{table}[t]
  \centering
  \begin{tabular}{l}
    \toprule
    Length tuple (hours) \\
    \midrule
    $(1, 6, 24)$ \\
    $(1, 3, 6, 12, 24)$ \\
    $(1, 6, 24, 48, 168)$ \\
    $(1, 2, 4, 8, 16, 24, 48, 96, 168)$ \\
    \bottomrule
  \end{tabular}
  \caption{Window length tuples used in the design grid.}
  \label{tab:length_sets}
\end{table}

\subsection{Window shapes}
Window shapes define how the history for an impression at hour $H$ is converted into time windows.
For a window of length $L$ hours, a trailing window uses the half open interval $[H-L, H)$.

Table \ref{tab:window_shapes} summarizes the benchmarked shapes.
Shapes that ``add'' windows are applied on top of trailing windows for the same length tuple.

\begin{table}[t]
  \centering
  \begin{tabular}{p{0.18\linewidth}p{0.76\linewidth}}
    \toprule
    Shape & Definition used in this study \\
    \midrule
    trailing &
      For each $L$ in the length tuple, use events in $[H-L, H)$. \\
    gap1 &
      For each $L$ in the length tuple, use events in $[H-(L+1), H-1)$.
      This excludes the most recent hour. \\
    bucket &
      Let $0 < e_1 < \dots < e_K$ be the hour edges from the length tuple.
      Construct disjoint buckets $[H-e_1, H)$, $[H-e_2, H-e_1)$, \dots, $[H-e_K, H-e_{K-1})$.
      Features are computed per bucket. \\
    calendar &
      Add calendar aligned day windows on top of trailing windows.
      For day based lengths, such as $L \in \{24, 48, 168\}$, the window is aligned to full days strictly before the impression hour.
      For sub day lengths, the calendar window matches the trailing window. \\
    event50 &
      Add an event count window on top of trailing windows.
      For each entity, use the most recent 50 impressions strictly before hour $H$. \\
    \bottomrule
  \end{tabular}
  \caption{Window shape definitions.}
  \label{tab:window_shapes}
\end{table}
